\section{Introduction}\label{sec:intro}

\paragraph{Main result.} The note studies pure Yang--Mills theory on a four-dimensional lattice torus, with the Wilson action $S$, by Ba\l aban's renormalization group with his symmetric block averaging \cite{BalYM4I,BalYM4II}. It treats the small-field part of this renormalization group. Its results concern the \emph{small-field model}, a modification of the lattice theory in which every step keeps only small fields: the curvature of the coarse configuration is at most a threshold $h_k$, and the fluctuation field and the white noise of the step are bounded by cutoffs. The densities of this model are
\[
\rho_0=\exp\big(-g_0^{-2}S-E_0\big),\qquad\rho_{k+1}=\one_{\mathfrak U_{k+1}(h_{k+1})}\,T^{\rm gf}\big(\chi^{\rm w}\chi_\xi\,\rho_k\big),\quad k=0,1,\dots,K-1,
\]
with the bare coupling $g_0$, a constant $E_0$, the renormalization group map $T^{\rm gf}$ with Ba\l aban's exponential gauge fixing, and the cutoffs $\chi^{\rm w}$ on the fluctuation field and $\chi_\xi$ on the white noise (Definition~\ref{def:small-field-step}, Remark~\ref{rem:sf-reference}(c)). What a step discards is nonnegative, so the partition function of the small-field model is at most that of the lattice theory with the same $E_0$ (Remark~\ref{rem:sf-reference}(c)). For a compact, connected and simple gauge group the note proves two results about this model.
\begin{itemize}
\item \emph{The flow} (Theorem~\ref{thm:N}). The couplings run by $g_{k+1}^{-2}=g_k^{-2}-\beta_{k+1}$, where $\beta_{k+1}$ is read off from the terms that the step creates. As long as the couplings stay in $(0,\gamma]$, with $\gamma$ small as in Remark~\ref{rem:flow-constants}, the density of every scale has Ba\l aban's generation form (Definition~\ref{def:gen-form}), each generation is bounded at the scale where it is created, and the $\beta_k$ are bounded. For the counterterm $E_0$ that makes the vacuum energy of the last scale $K$ vanish, the partition function is an explicit product of normalizations times a number between $e^{-c|\T^{(K)}|}$ and $e^{c|\T^{(K)}|}$, where $|\T^{(K)}|$ is the number of sites of the last lattice. The constants do not depend on the scale, the number of steps, the torus or the couplings, except that $c$ depends on the last coupling $g_K$ (Theorem~\ref{thm:N}(vi)).
\item \emph{The ultraviolet behaviour of the couplings} (Theorem~\ref{thm:M}(iv)). Let $p_0(g)=(\log g^{-1})^q$ be the white-noise cutoff. For every small final coupling $g$ and every number $K$ of steps there is a bare coupling $g_0$ for which the flow stays in $(0,\gamma]$ and ends at $g_K=g$, and
\[
\frac1{g_k^2}=\frac1{g^2}+(K-k)\beta_\infty+O\Big(\big(\log(K-k+2)\big)^{6q+1}\Big)\quad(k\le K),\qquad\beta_\infty=\frac{11C_A}{24\pi^2T_\rho}\log L,
\]
with a constant independent of $K$ and of the torus. Here $C_A$ is the Casimir of the adjoint representation, and $T_\rho$ is the normalization $\operatorname{tr}_\rho(T^aT^b)=T_\rho\delta^{ab}$ of the trace of the representation $\rho$ in the action (Theorem~\ref{thm:M}). Here $\beta_\infty$ is the one-loop coefficient per step (Theorem~\ref{thm:M}(iii)). The coupling $g$ of the action $S$ is related to the standard one by $g_{\rm std}^2=g^2/T_\rho$ (Section~\ref{sec:ol-feynman}, after Theorem~\ref{thm:ol-log}; Section~\ref{g:result}).
\end{itemize}
The torus is fixed at the last scale, with a side of at least $n_{\min}$ there (Section~\ref{subsec:tf-scales}), and all bounds are uniform in the number of scales (Remark~\ref{rem:sf-flow}(d)). With $K=N$ steps on a torus of physical side $\ell$, the finest spacing is $L^{-N}$ (Section~\ref{subsec:tf-scales}); so the bounds are uniform in the lattice spacing at fixed physical volume. Every hypothesis of Theorem~\ref{thm:L} holds at every step of the flow (Theorem~\ref{thm:N}(v)), and the hypothesis $g_k\le\gamma$ of Theorem~\ref{thm:N} holds for the bare couplings of Theorem~\ref{thm:M}(iv).

\emph{What is not proved.} The large-field stage is not treated, so nothing is proved about the partition function or the correlation functions of the lattice theory itself (Remark~\ref{rem:sf-reference}(d)). No limit of the measure or of correlation functions as the lattice spacing tends to $0$ is considered: Theorem~\ref{thm:M}(iv) concerns the couplings of the small-field flow (Section~\ref{ss:1.5}).

\emph{How to read.} The \hyperref[guide]{guide to the proof} follows this introduction. Section~\ref{g:read} gives four reading paths, Section~\ref{g:result} states the result in more detail, Section~\ref{g:scales} the scales, and Section~\ref{g:strategy} the strategy of the four-dimensional flow. Sections~\ref{g:linear}--\ref{g:oneloop} follow the Parts, and Sections~\ref{g:parameters}--\ref{g:hard} give the order of the parameters, the notation (with Appendix~\ref{app:notation}), the comparison with Ba\l aban and Dimock, and the places where the work is hardest. The companion volume \emph{Details} expands the proofs of Parts~VII--XI step by step (Section~\ref{g:read}).

\subsection{Ba\l aban's renormalization group}\label{ss:1.1}

Ba\l aban's construction of lattice Yang--Mills theory \cite{BalPropI,BalPropII,BalAvg,BalReg,BalBackground,BalVar,BalUV3,BalYM4I,BalYM4II} is a renormalization group (RG) analysis. A gauge field on a fine lattice of spacing $\eta=L^{-k}$ is averaged over blocks of side~$1$, and this averaging is performed in $k$ steps, each enlarging the blocks by a factor~$L$. After $k$ steps one has to integrate over all fine configurations whose block averages equal a given coarse configuration. Following the usual Gaussian strategy, one expands the action around its minimizer subject to the averaging constraint. This minimizer is called the \emph{background field} \cite{BalVar}. The fluctuations around it are then Gaussian to leading order, with an explicit covariance.

Every estimate in this programme rests on properties of a small family of \emph{linear} operators:
\begin{itemize}
\item Green functions of the form $(\text{Laplace-type operator}+a\,Q^*Q)^{-1}$, where $Q$ is the $k$-step averaging operator;
\item the \emph{minimizer map} $H$, which sends a coarse configuration $B$ to the fine configuration of least action with average~$B$;
\item the \emph{constrained Green function}, i.e.\ the covariance of the fluctuation field, which has vanishing averages;
\item inverses of averaged Green functions, such as $(QGQ^*)^{-1}$, which live on the coarse lattice.
\end{itemize}
The properties needed are \emph{exponential decay} on the unit scale and \emph{local regularity}: bounds on values, first derivatives and H\"older quotients of first derivatives. They must be uniform in the number $k$ of RG steps, which means uniform in the ratio $\eta^{-1}$ between the unit scale and the lattice scale, and uniform in the size of the volume.

For the full non-abelian problem these operators depend on the background field $U$. Ba\l aban obtained them by perturbation from the case $U=1$ \cite{BalPropI}, through random-walk expansions \cite{BalPropII,BalBackground}. The case $U=1$ is the foundation of the whole analysis. It is also the case in which all the scales are already present, together with the gauge structure that distinguishes gauge fields from scalar fields.

In four dimensions the Wilson action carries no power of the lattice spacing, so a change of units does not change the coupling: the coupling is marginal, and the coupling of the next scale is produced by the fluctuation integral alone (Section~\ref{sec:rg-step}). The terms that a step creates are of order one, and the iteration does not contract them. Ba\l aban keeps them as generations, each bounded at the scale where it is created, and cancels their marginal parts at the later scales by counterterms proportional to the action \cite{BalYM4I,BalYM4II} (Remark~\ref{rem:flow-status}). Every step has a small-field part, in which the fields are close to the background field, and a large-field part. The note treats the small-field part.

\subsection{What this note does}\label{ss:1.2}

The note carries out the small-field part of Ba\l aban's analysis of four-dimensional lattice Yang--Mills theory on a single torus, without the nested domains $\Omega_j$ of \cite{BalVar}: from the linear operators to the flow over all scales and the one-loop coefficient. It has eleven Parts, in four layers.
\begin{itemize}
\item \emph{The linear operators} (Parts~I--III, Theorems~\ref{thm:A}--\ref{thm:G}): Green functions with an averaging term, minimizer maps and constrained Green functions, with exponential decay and pointwise bounds uniform in the number of averaging steps and in the torus; at $U=1$, at backgrounds close to $1$ in a global gauge, and at locally regular backgrounds.
\item \emph{The background field} (Parts~IV--VI, Theorems~\ref{thm:H}--\ref{thm:J}): Ba\l aban's symmetric nonlinear averaging, the axial-to-Landau chart, and the variational problem. For a real coarse configuration of small curvature the minimizers of the action with prescribed averages form exactly one critical orbit among the configurations of small curvature and current, and they are locally regular, uniformly in the number of steps.
\item \emph{One step and its localization} (Parts~VII--IX, Theorem~\ref{thm:K}): one step of the renormalization group as an exact Gaussian integral around the background field, with a coercive fluctuation operator and an exponentially decaying covariance; the localization of its dependence on the background, of radius type (Part~VIII) and of volume type, by random walks with weakening parameters (Part~IX).
\item \emph{The flow and the one-loop coefficient} (Parts~X--XI, Theorems~\ref{thm:L}, \ref{thm:N} and~\ref{thm:M}): the step in polymer form, its iteration over all scales in Ba\l aban's generation form, the one-loop coefficient of the running coupling, and the ultraviolet behaviour of the couplings.
\end{itemize}
Figure~\ref{fig:G-mainline} shows how these results feed Theorems~\ref{thm:N} and~\ref{thm:M}(iv), and Table~\ref{tab:G-parts} lists the main results of each Part.

\paragraph{Scope.} The dimension is $d=4$ throughout, and the gauge group $G$ is compact and connected; from Part~VII on it is semisimple, and for the flow of Part~X (Sections~\ref{sec:generations}--\ref{sec:flow4}) and in Part~XI simple. Semisimplicity gives the quadratic bound for normalized activities (Lemma~\ref{lem:quadratic-bound}). For simple $G$ a quadratic form in the curvature that is invariant under conjugation by $G$ and under $B_4$ is a multiple of a single isotropic form (Lemma~\ref{lem:b4-forms}(a)); so the period density of a generation is a multiple of that of $S$, and the running coupling is one number (Proposition~\ref{prop:period-iso}). The flow results of Parts~VII--X and Theorem~\ref{thm:M}(iv) concern the small-field model on a fixed torus; Part~XI also uses tori of growing side, through the torus independence of Proposition~\ref{prop:one-loop-local} (Part~IX), and infinite-volume limits of the one-loop terms (Theorem~\ref{thm:M}(ii),~(iii)). The large-field stage is not treated (Section~\ref{ss:1.5}).

\paragraph{The trivial background (Part~I).} Part~I treats the linear operators at the trivial background $U=1$, on a single torus with $k$-fold averaging everywhere. It is self-contained, and it is meant to be readable by a PhD student who knows linear algebra, Fourier series and the basics of elliptic equations. All lattice computations are written out.

Let us make precise what ``$U=1$'' means. The gauge field is $U=e^{\ii\eta A}$ with $A$ taking values in the Lie algebra $\mathfrak g$ of the gauge group. Expanding the Wilson action to second order around $U=1$ gives $\tfrac12\|dA\|^2$, and infinitesimal gauge transformations act as $A\mapsto A+d\lambda$. The non-abelian structure enters only at higher orders. At $U=1$ every operator acts componentwise in a basis of $\mathfrak g$. Hence it suffices to treat real-valued fields, and this is what we do. This is \emph{not} abelian gauge theory. It is the linearization of non-abelian theory at its trivial background. The operators obtained are those which Ba\l aban perturbs around in \cite{BalBackground,BalVar}.

\paragraph{The operators of Part~I.} The precise definitions are given in Sections~\ref{sec:setup}--\ref{sec:gauge-coercivity}. Here is an overview. We use a fine torus $\Tf$ of spacing $\eta$ and a unit torus $\T\subset\Tf$, the scalar block average $Q_0$, and the gauge-field average $\bQ$ of \cite{BalPropI}. The averaged gauge field $(\bQ A)_\mu$ is an average of $A_\mu$ over the block and along lines of unit length in direction~$\mu$. It satisfies $\bQ\,d=\dl Q_0$, where $\dl$ is the unit-lattice gradient (gauge covariance). Since the blocks are centred, $Q_0$ and $\bQ$ commute with the translations of $\T$ and with the hyperoctahedral group $B_4$ acting about the block centres (Lemma~\ref{lem:avg}(f)). We study:
\begin{enumerate}[label=(\roman*)]
\item the scalar Green functions $G_\sigma=(-\Lap+aQ_\sigma^*Q_\sigma)^{-1}$. Here $Q_\sigma$ is either the block average $Q_0$ or one of the ``tube'' averages $Q_1,\dots,Q_d$ that make up $\bQ$;
\item the operators derived from them: $(Q_\sigma G_\sigma Q_\sigma^*)^{-1}$, the minimizer map $H_\sigma$, the constrained Green function $C_\sigma$, and the effective action. We also study the projection $R$ which defines Ba\l aban's Landau-type gauge;
\item the gauge-field operator $\bL_a=d^*d+dRd^*+a\bQ^*\bQ$ and its inverse $\bG$;
\item the gauge-field minimizer map $\bH=\bG\bQ^*(\bQ\bG\bQ^*)^{-1}$, the effective action $\bD$, and the constrained Green function $\bC$, which is the operator appearing in the Euler--Lagrange equation of the variational problem \cite[(110)--(111)]{BalVar}.
\end{enumerate}

\subsection{Main results}\label{ss:1.3}

The theorems of Parts~I and~II are stated with the conventions of Part~I: $d=4$, $L\ge3$ is odd, $k\ge1$, $\eta=L^{-k}$, and the unit torus has side $n\ge8$. Since $L$ is odd, $N=L^k$ is odd. The parameter $a>0$ in front of the averaging term is fixed. Ba\l aban takes $a=1$. Every constant in Part~I depends only on $a$ and a H\"older exponent $\alpha\in(0,1)$; constants that depend only on the dimension are absolute (Section~\ref{sec:setup}). In Part~II it depends in addition on $\dim\mathfrak g$ and on the decay rate $\gamma_W$ of the potential term $W$ of Theorem~\ref{thm:F}. In particular, the constants of Parts~I and~II do not depend on $L$, $k$ or $n$. The constants of Theorems~\ref{thm:reg-scalar}, \ref{thm:reg-coercive} and~\ref{thm:reg-gauge} (on the torus) and~\ref{thm:reg-min} of Part~III do not depend on $L$, $k$, $n$ or $M\ge M_0$ either (Lemma~\ref{lem:uniformity}(a)); the random-walk expansion and Theorem~\ref{thm:reg-locality}(b) have constants that depend on $M$ (Remark~\ref{rem:uniformity-limits}(b)). From Part~IV on, constants may depend on $L$ (Section~\ref{sec:averaging}), and the flow needs $L$ large (Remark~\ref{rem:parameters}(2)); the order in which all parameters are chosen is in Section~\ref{g:parameters}.

For a fine-lattice field $u$ (scalar or vector-valued) and a set $X$ we write
\[
|u|_X=\sup_{x\in X}|u(x)|,\qquad
|\partial u|_X=\max_\mu\sup_{x\in X}|\partial_\mu u(x)|,\qquad
\Hol{v}{X}=\sup_{x\ne x'\in X}\frac{|v(x)-v(x')|}{|x-x'|^\alpha}.
\]
Here $\partial_\mu$ is the forward lattice derivative, and the maximum also runs over components. $B(y)$ is the unit block centred at $y\in\T$, the cube of $N^d$ fine points with centre $y$, and $|y-y'|$ is the distance on the torus. For a gauge field $B$ on $\T$, $|B|_{y'}$ is the maximum of $|B|$ over the unit bonds starting at $y'$.

\begin{maintheorem}[Coercivity; Corollary~\ref{cor:scalar-coercive}, Theorem~\ref{thm:gauge-coercive}]\label{thm:A}
There is $c>0$ such that for all scalar fields $f$ and all $\sigma\in\{0,\dots,d\}$,
\[
\|\partial f\|^2+a\|Q_\sigma f\|^2\ \ge\ c\,\big(\|f\|^2+\|\partial f\|^2\big),
\]
and for all gauge fields $A$,
\[
\ip{A}{\bL_a A}=\|dA\|^2+\|Rd^*A\|^2+a\|\bQ A\|^2\ \ge\ c\,\big(\|A\|^2+\|\partial A\|^2\big).
\]
\end{maintheorem}

\begin{maintheorem}[Scalar Green functions; Theorem~\ref{thm:scalar-green}]\label{thm:B}
There are $C,\gamma>0$ such that for all $y,y'\in\T$ and all $f$ supported in $B(y')$:
\begin{align*}
|G_\sigma f|_{B(y)}+|\partial G_\sigma f|_{B(y)}+\Hol{\partial G_\sigma f}{B(y)}&\le Ce^{-\gamma|y-y'|}\,|f|_{B(y')},\\
|G_\sigma \partial_\mu^* f|_{B(y)}+\Hol{G_\sigma\partial^*_\mu f}{B(y)}&\le Ce^{-\gamma|y-y'|}\,|f|_{B(y')}.
\end{align*}
Corresponding $L^2$ bounds hold for $G_\sigma$, $\partial G_\sigma$, $G_\sigma\partial^*$ and $\partial G_\sigma\partial^*$ between blocks.
\end{maintheorem}

\begin{maintheorem}[Averaged scalar operators; Section~\ref{sec:scalar-avg}]\label{thm:C}
The operator $Q_\sigma G_\sigma Q_\sigma^*$ on the unit torus satisfies $c\le Q_\sigma G_\sigma Q_\sigma^*\le C$. Its inverse has a kernel bounded by $Ce^{-\gamma|y-y'|}$. The minimizer map $H_\sigma$ satisfies
\[
|H_\sigma g|_{B(y)}+|\partial H_\sigma g|_{B(y)}+\Hol{\partial H_\sigma g}{B(y)}\le C\sum_{y'}e^{-\gamma|y-y'|}|g(y')|,
\]
and the derivative bounds hold with $|g(y')|$ replaced by the unit-lattice gradient $|\dl g(y')|$. The constrained Green function $C_\sigma$ satisfies the bounds of Theorem~\ref{thm:B}. The effective action $\cD_\sigma$, defined by $\ip{g}{\cD_\sigma g}=\min\{\|\partial f\|^2:Q_\sigma f=g\}$, satisfies $c\|\dl g\|^2\le\ip{g}{\cD_\sigma g}\le C\|\dl g\|^2$ and has an exponentially decaying kernel. The projection $R$ is $I$ minus an operator with exponential decay.
\end{maintheorem}

\begin{maintheorem}[Gauge-field Green function; Theorem~\ref{thm:gauge-green}]\label{thm:D}
The operator $\bG=\bL_a^{-1}$ satisfies, for $J$ supported in $B(y')$,
\[
|\bG J|_{B(y)}+|\partial\bG J|_{B(y)}+\Hol{\partial\bG J}{B(y)}\le Ce^{-\gamma|y-y'|}\,|J|_{B(y')},
\]
together with $L^2$ block bounds for $\bG$, $\partial\bG$, $\bG\partial^*$ and $\partial\bG\partial^*$.
\end{maintheorem}

\begin{maintheorem}[Gauge-field minimizer and fluctuation operators; Section~\ref{sec:gauge-min}]\label{thm:E}
One has $c\le\bQ\bG\bQ^*\le C$, and $(\bQ\bG\bQ^*)^{-1}$ has an exponentially decaying kernel. The minimizer map $\bH$ sends $B$ to the minimizer of $\|dA\|^2$ subject to $\bQ A=B$ and $Rd^*A=0$, and it satisfies
\[
|\bH B|_{B(y)}+|\partial\bH B|_{B(y)}+\Hol{\partial\bH B}{B(y)}\le C\sum_{y'}e^{-\gamma|y-y'|}|B|_{y'}.
\]
The gauge-field effective action satisfies $c\|\dl B\|^2\le\ip{B}{\bD B}=\|d\bH B\|^2\le C\|\dl B\|^2$. The constrained Green function $\bC=\bG-\bG dRd^*\bG-\bH\bQ\bG$ is the Green function of $\|dA\|^2$ on $\{\bQ A=0,\ Rd^*A=0\}$, and, for $J$ supported in $B(y')$,
\[
|\bC J|_{B(y)}+|\partial\bC J|_{B(y)}+\Hol{\bC J}{B(y)}\le Ce^{-\gamma|y-y'|}\,|J|_{B(y')}.
\]
\end{maintheorem}

These are, at $U=1$ and on a single torus, the bounds that Ba\l aban states in \cite[Sect.~F]{BalPropI} and \cite[Theorems~3.12, 3.13]{BalBackground} and uses in \cite{BalVar}.

\paragraph{Backgrounds (Part II).} In Part~II the fields take values in the complexified Lie algebra $\mathfrak g^c$. The background $U=e^{\ii\eta A}$ satisfies $|A|,|\partial A|,[\partial A]_\alpha\le\eps$ in a global gauge (hypothesis (B)). Hypothesis (B) is restrictive. It forces all Wilson loops to obey a perimeter bound, so the background must be close to flat on scales larger than $O(1)$ (Section~\ref{sec:background}). The operators now involve parallel transport. Part~II is stated for a model averaging pair $(Q'(U),\bQ(U))$, with transports along a single coordinate-ordered contour from the block centre and only the main term of Ba\l aban's linearized averaging. It transforms exactly under gauge transformations, but it intertwines the covariant derivatives only up to a defect: $\bQ(U)D_U=D_VQ'(U)+E(U)$ with $E(U)=O(\eps)$. For the linearization of Ba\l aban's symmetric averaging $E=0$ (Part~IV), and the results apply to it (Remark~\ref{rem:general-averaging}): the real statements under properties (i)--(iv) of that remark, the analytic ones with property (v) in addition.

\begin{maintheorem}[Backgrounds; Sections~\ref{sec:cov-scalar}--\ref{sec:cov-min}]\label{thm:F}
Let $U$ satisfy (B) with parameter $\eps$, and let $W\in\cW(\eps,\gamma_W)$ in this gauge (Section~\ref{sec:background}). For $\eps$ small enough, with constants independent of $L,k,n$ and $U$:
\begin{enumerate}[label=(\roman*)]
\item $G'(U)=(D^TD+aQ'^TQ')^{-1}$ satisfies the bounds of Theorem~\ref{thm:B}. The inverse $\Gamma'(U)$ of $Q'G'Q'^T$, the minimizer map $H'(U)$, the constrained Green function $C'(U)$ and the projection $R(U)$ satisfy the decay and pointwise bounds of Theorem~\ref{thm:C} (Theorem~\ref{thm:cov-scalar}, Proposition~\ref{prop:cov-R}). The bounds in terms of $|\dl g|$ and the effective-action bounds of Theorem~\ref{thm:C} are stated only at $U=1$.
\item The gauge-field Green function $\bG(U)$ satisfies the bounds of Theorem~\ref{thm:D} (Theorem~\ref{thm:cov-gauge}).
\item The inverse of $\bQ_\pi\bG\bQ_\pi^T$, the minimizer map $\bH(U)$ and the constrained Green function $\bC(U)$ satisfy the decay and pointwise bounds of Theorem~\ref{thm:E} (Theorem~\ref{thm:cov-min}). For every $J$, $\bC(U)J$ is the unique $A\in\cX(U)$ with $q_U(A',A)=\mathrm b(A',J)$ for all $A'\in\cX(U)$. For real $U$, real symmetric $W$ and real data $B$, $\bH(U)B$ is the unique minimizer of $q_U(A,A)$ among real $A$ with $\bQ(U)A=B$ and $R(U)D_U^TA=0$. The effective-action bounds of Theorem~\ref{thm:E} are stated only at $U=1$.
\end{enumerate}
Here $\mathrm b(A,A')=\eta^d\sum_{x,\mu}\operatorname{tr}\big(A_\mu(x)A'_\mu(x)\big)$ is the bilinear form, $q_U(A,A')=\mathrm b(D_UA,D_UA')+\mathrm b(A,WA')$ is the quadratic form, with the small potential term $W$, $\cX(U)=\{A:\bQ(U)A=0,\ R(U)D_U^TA=0\}$, and $\bQ_\pi=\bQ(U)\Pi(U)$. For real $U$, real symmetric $W$ and real $A$, $q_U(A,A)=\|D_UA\|^2+\ip A{WA}$. The projection $\Pi(U)$ onto Ba\l aban's gauge slice $\{R(U)D_U^TA=0\}$ along the residual gauge directions is Ba\l aban's gauge-invariant extension \cite[(3.118)--(3.119)]{BalBackground}. All these operators depend analytically on $A$ and $W$ and transform covariantly under $G$-valued gauge transformations, with $W$ transformed accordingly. The identities $\bG DR=DG'R$, $RD^T\bG=RG'D^T$ and $\bQ_\pi\bG DR=0$ hold exactly.
\end{maintheorem}

The operators are covariant under $G$-valued gauge transformations $u$, with $W$ replaced by $R(u)WR(u)^{-1}$; the latter need not belong to $\cW(\eps,\gamma_W)$ in the new gauge. Hence Theorem~\ref{thm:F} transfers to every background gauge equivalent to one satisfying (B), as follows (Proposition~\ref{prop:covariance}): the bounds on values and the $L^2$ block bounds hold with $\partial$ replaced by the covariant derivative, the identities hold, and the H\"older bounds hold in the gauge in which (B) holds. The comparisons with the operators at $U=1$ proved in Sections~\ref{sec:cov-scalar}--\ref{sec:cov-min} hold only in that gauge; they fail for a pure gauge $U=1^u$.

The proofs are perturbative: every operator is written as its $U=1$ counterpart plus a perturbation of a specific structure (Section~\ref{sec:perturbation}), and three tools of Part~I, the Combes--Thomas method, inversion with decay and local regularity, are shown to be stable under such perturbations. For the pointwise bounds a further regularity estimate is needed, a lattice Schauder estimate for divergence-form equations (Theorem~\ref{thm:local-regularity}(d)).

\paragraph{Local regularity (Part III).} In Part~III the global hypothesis (B) is replaced by a local one. A real background is \emph{$(M,\eps)$-regular} if on every cube of side of order $M$ there is a gauge in which (B) holds; the gauges on different cubes are unrelated. This is the form of Ba\l aban's regularity conditions on the background.

\begin{maintheorem}[Regular backgrounds; Sections~\ref{sec:reg-scalar}--\ref{sec:random-walk} and~\ref{sec:nested-walks}]\label{thm:G}
There is $M_0$ such that for $M\ge M_0$ and $(M,\eps)$-regular backgrounds with $\eps$ small (independent of $M$) the following hold.
\begin{enumerate}[label=(\roman*)]
\item The statements of Theorem~\ref{thm:F} remain true, with local potentials $W\in\cW_{\rm loc}(\eps)$ (Section~\ref{sec:reg-gauge}). The comparisons with the operators at $U=1$ of Sections~\ref{sec:cov-scalar}--\ref{sec:cov-min} need a global gauge and fail for a pure gauge $U=1^u$; they are stated only under~(B), in those sections. The scalar operators need only $(1,\eps)$-regularity. Sup norms of derivatives are taken with the covariant derivative, and H\"older quotients in local gauges.
\item The operators whose inverses are $G'(U)$ and $\bG(U)$ have Dirichlet versions on arbitrary unions $X$ of blocks. These depend only on the background on the bonds of $\hat X$, the set of points at $\infty$-distance at most $1$ from $X$, and they are coercive (the gauge-field version with a mass term near the boundary and Dirichlet conditions for all components of the gauge field). Their inverses, the inverse of the averaged scalar Green function on $X$, and $\one_X-R_X$, with $R_X$ the Dirichlet version of the projection $R$, have exponential decay uniformly in the domain.
\item \emph{Locality in the background:} if two regular backgrounds agree on a region, the propagators, the minimizer map and the fluctuation operator agree deep inside it up to errors that are exponentially small in the distance to the complement.
\item Ba\l aban's random-walk expansion converges for $M$ large: for the scalar Green function $G'(U)$ as a sum of products of local propagators (Theorem~\ref{thm:random-walk}(a)); for the gauge-field Green function $\bG(U)$ in the factorized form of Section~\ref{sec:nested-walks}, with the domains $\Lam^v_j$ of Section~\ref{sec:weakenable} for a fixed $\rho\ge15$, so that \eqref{eq:rho-choice} holds with $w=10M$, for $M$ large depending on $\rho$, on tori of side $n\ge C(\rho+C_v)M$, with $C_v$ the margin of level~(L3) there (Theorem~\ref{thm:random-walk}(b)). Under the same hypotheses, $\Gamma'(U)$, $H'(U)$ and $R(U)$ are expanded at levels~(L1) and~(L2) there. Section~\ref{sec:nested-walks} is written for $E=0$; the additional terms for $E\ne0$ are treated in the proof of Theorem~\ref{thm:random-walk}(b).
\item All results extend analytically to complex perturbations $Ue^{\ii\eta\delta A}$ of a regular background by a field $\delta A$ that is small in the local gauges, for potentials $W$ that depend analytically on the background, with coercivity replaced by $\Rea$-coercivity.
\end{enumerate}
\end{maintheorem}

The key new step is a local proof of coercivity (Section~\ref{sec:reg-gauge}). An IMS localization at scale $M$ reduces it to cubes. On each cube one works in the local gauge, perturbs around $U=1$, and uses the coercivity of Part~I. Fourier analysis is used only at $U=1$.

\paragraph{The nonlinear problem (Parts IV and V).} Ba\l aban's averaging of a gauge field is nonlinear and is performed one scale at a time. We use his symmetric four-dimensional averaging \cite[(0.3)--(0.12)]{BalYM4I} (Section~\ref{sec:averaging}). Its blocks are centred at the coarse points. The transport from a block centre $y$ to a point $x$ of the block is the group average (Lemma~\ref{lem:group-average}) of the transports along the coordinate-ordered contours from $y$ to $x$, one for each order of the directions (Definition~\ref{def:families}), and the coarse bond variables are built from these averaged contour variables. The $k$-fold composition $U\mapsto U^{(k)}$ is covariant with respect to the values of gauge transformations at the points of $\T$, and under the group $\mathfrak S_k$ generated by the translations of $\T$ and by $B_4$ about its points (Lemma~\ref{lem:one-step-basic}). So the residual gauge group of the variational problem is $\mathfrak G_0=\{u:u|_\T=1\}$. Part~IV constructs the bridge to the linear theory. Following \cite{BalAvg,BalReg} and \cite[pp.~253--254]{BalYM4I}, we introduce the averaged axial gauge, defined through the averaged contour variables, and the gauge-adapted averaging $\mathcal Q$. We show (Proposition~\ref{prop:linearized-averages}, Corollary~\ref{cor:parts-apply}) that the linearization of $\mathcal Q$ at a regular reference $U_0$ is an averaging pair $(Q'_k(U_0),\bQ_k(U_0))$ to which Parts~II and~III apply, with \emph{exact} covariance ($E=0$), and that it is covariant under $\mathfrak S_k$. We then prove Ba\l aban's axial-to-Landau theorem \cite[Theorem~2]{BalReg} on a single torus.

\begin{maintheorem}[Axial-to-Landau chart; Theorem~\ref{thm:chart}]\label{thm:H}
Let $U_0$ be $(M,\eps)$-regular with current at most $\eps$, where $\eps$ is small. Every real configuration $U$ with small curvature and current, and with $k$-fold averages close to those of $U_0$, has exactly one small representative $e^{\ii\eta A}$ relative to $U_0$ that lies in Ba\l aban's Landau gauge $R(U_0)D^T_{U_0}A=0$ and is related to the averaged axial representative of $U$ by the correspondence transformation. It obeys $|A|+|\nabla A|+|\Lap_{U_0}A|+[\nabla A]_\alpha\le C(\eps+\text{curvature}+\text{current}+\text{difference of averages})$, uniformly in $k$, $n$ and $\eta$. The curvature and the current in this bound are those of $U$.
\end{maintheorem}

\begin{maintheorem}[The variational problem relative to a regular reference; Theorem~\ref{thm:variational}]\label{thm:I}
Let $U_0$ be $(M,\eps)$-regular with current at most $\eps$, and let $V$ be a real configuration on $\T$ with $V(U_0^{(k)})^{-1}=e^{\ii B}$, $|B|\le\beta$. If $\eps+\beta$ is small, then the Wilson action restricted to $\{U:U^{(k)}=V\}$ has exactly one critical $\mathfrak G_0$-orbit among configurations of small curvature and current. The orbit consists of the minimizers of the action among these configurations. Its Landau representative relative to $U_0$ satisfies
\[
|A_*|+|\nabla A_*|+|\Lap_{U_0}A_*|+[\nabla A_*]_\alpha\le C(\eps+\beta),\qquad A_*=\bH B-\bC J_{U_0}+O\big((\eps+\beta)^2\big).
\]
The minimizer is again regular, it depends analytically on $B$, and the derivative with respect to $B$ decays exponentially. The critical orbit, the Landau representative and the multiplier are covariant under $\mathfrak S_k$.
\end{maintheorem}

This is the content of Ba\l aban's variational theorem \cite[Theorem~1]{BalVar} for one step $k$ of his induction, given the reference configuration $U_0$, on a single torus. Along the way we reprove Ba\l aban's results on the averaging \cite{BalAvg} and on gauge fixing \cite{BalReg} in the form needed.

\paragraph{The induction (Part VI).} Ba\l aban proves his theorem by induction over $k$. The reference at step $k$ is the solution of the $(k-1)$-step problem for the trivial extension of $V$; it has exactly the prescribed averages. Theorem~\ref{thm:I} then gives the solution at step $k$. Its regularity bound is worse by a power of $L$, because units change by a factor $L$. Ba\l aban restores the bound by a local analysis on nested domains \cite[Sect.~F]{BalVar}. We stay on the torus and use the uniformity of Theorem~\ref{thm:I} in the reference. Near a given point we replace the reference by a configuration that is trivial, in a local gauge, near the point and equals the solution far away. Relative to it, the solution near the point is determined by the coarse field in the gauge of the coordinate-ordered contours from the point, which is controlled by the curvature of $V$ alone; this gauge is a device of the proof (Section~\ref{sec:local-gauges}). The errors are exponentially small in the splicing radius, or quadratic.

\begin{maintheorem}[Ba\l aban's variational theorem; Theorem~\ref{thm:balaban}, Corollary~\ref{cor:regularity}]\label{thm:J}
There are $a_1,B_3,n_0$ such that the following holds on tori of side $n\ge n_0$, for every $k$ and every real configuration $V$ on $\T$ whose curvature is at most $\eps_1\le a_1$. The Wilson action on $\{U:U^{(k)}=V\}$ has exactly one critical $\mathfrak G_0$-orbit among configurations of small curvature and current. It consists of the minimizers of the action among these configurations, which have curvature and current at most $B_3\eps_1$ and are $(M,B_3\eps_1)$-regular. For every $r\in\mathfrak S_k$ the minimal orbit of $rV$ is the image under $r$ of the minimal orbit of $V$. The constants do not depend on $k$, $n$ or $\eta$. Moreover, every real configuration whose curvature and current are at most $\theta$, with $\theta$ small, is $(M,C\theta)$-regular.
\end{maintheorem}

The last statement is the torus version of Ba\l aban's derivation of local regularity from gauge-invariant conditions \cite[Sect.~F]{BalReg}. So the hypothesis of Part~III holds for all configurations of Ba\l aban's small-field class.

\paragraph{One renormalization group step (Part VII).} The renormalization group transformation $T$ integrates over the configurations on $\T=\T^{(k)}$ with prescribed averages on $\T'=\T^{(k+1)}$ (Section~\ref{sec:rg-step}). In four dimensions the Wilson action carries no power of the lattice spacing, so a change of units does not change the coupling. The coupling of the next step is produced by the fluctuation integral, through the coefficient $\beta$ of the action in the effective density \cite[(1.20)--(1.22)]{BalYM4I}. The integrand is invariant under the gauge transformations $\mathcal G'$ of $\T$ that are trivial on $\T'$. Following Ba\l aban \cite[(0.14)--(0.19)]{BalYM4I}, we fix this invariance by an exponential gauge-fixing function $\mathrm{gf}_W$: a product over the blocks of weights of the form $e^{-|u-1|^2/2g^2}$ in the averaged contour variables relative to the base configuration $W$ of the step (Definition~\ref{def:gauge-fixing}). By a Faddeev--Popov identity without determinant (Lemma~\ref{lem:exp-gauge}(a)), the gauge-fixed map $T^{\rm gf}$ agrees with $T$ on $\mathcal G'$-invariant integrands supported on regular configurations. The small-field step is $T_{\rm sf}(\rho)=T^{\rm gf}(\chi^{\rm w}\chi_\xi\rho)$, with bondwise cutoffs on the fluctuation field and on the white noise, as in \cite[(2.9), (2.10)]{BalYM4I} and \cite[(206)]{DimockI}. These cutoffs are not $\mathcal G'$-invariant, so the small-field steps define a modified model (Remark~\ref{rem:sf-reference}). We write the gauge-fixed step exactly as a Gaussian integral:
\begin{itemize}
\item the background of the step is the $(k+1)$-step background field $U'$ (Lemma~\ref{lem:composition});
\item the constraint is eliminated on the middle bond $b_0(c)$ of each coarse bond $c$, with local Jacobians (Lemma~\ref{lem:eliminate}). The fluctuation field lives on the set $\mathfrak F$ of the other bonds, which is invariant under the lattice symmetries (Definition~\ref{def:bonds});
\item the action becomes a function of the coarse fluctuation through the coarse effective action $V\mapsto S(U_k(V))$, whose gradient is the Lagrange multiplier and whose derivatives decay along trees (Proposition~\ref{prop:coarse-action}). The gauge fixing adds the quadratic form $\frac12\|\mathfrak G_WZ\|^2$, with the averaged transports $\mathfrak G_W$ of Definition~\ref{def:avg-transport}, and a remainder of third order that is a sum of terms localized in single blocks.
\end{itemize}

\begin{maintheorem}[One renormalization group step; Theorems~\ref{thm:step} and~\ref{thm:step-coercive}, Corollary~\ref{cor:covariance}, Lemma~\ref{lem:gf-step-invariance}, Proposition~\ref{prop:symmetry}]\label{thm:K}
Let $G$ be compact, connected and semisimple, let $M\ge\max(M_0,5L)$, and let the side of the torus satisfy the conditions of Section~\ref{sec:rg-step}. There are $a_{\rm step},r_2>0$, depending only on $L$, $G$, $a$, $\alpha$ and $M$, and $c_L>0$, depending only on $L$, such that the following holds for every $k\ge1$. Let $V'$ be a real configuration on $\T'$ with curvature at most $\eps'_1\le a_{\rm step}$, let $U'$ be its background field and $W$ the base configuration on $\T$, and let $w\le r_2$. Let $\rho$ be a gauge-invariant density of the form \eqref{eq:rho-form}, with coupling $g=g_k$ and threshold $\tilde h$, where $a_1\ge\tilde h\ge2(C_*+1)L^{-2}\eps'_1+4w$.
\begin{enumerate}[label=(\roman*)]
\item \emph{(Exact representation.)} $T^{\rm gf}(\chi^{\rm w}\rho)(V')$ is $\mathcal N_k\,e^{-g^{-2}S(U')}(\det\Delta_k)^{-1/2}\prod_c\mathfrak j_c$ times the integral, with respect to the Gaussian measure with covariance $C_k=\Delta_k^{-1}$ on the fields on $\mathfrak F$, of $\chi^{\rm w}$ times the exponential of an interaction of third order, of local Jacobian terms and of the old activities. The constant $\mathcal N_k$ is explicit and does not depend on the background; the $\mathfrak j_c$ are local Jacobians, and $\log\det\Delta_k$ is a sum of quasi-local, gauge-invariant densities, which depend on the whole background with exponential decay.
\item \emph{(Coercivity.)} $\frac14c_L\|Z\|^2\le\Delta_k(Z,Z)\le C_L\|Z\|^2$ for real fields $Z$ on $\mathfrak F$, and $C_k$ has exponential decay; at complex bases close to $W$, $\Rea\Delta_k\ge\frac18c_L$. The constant $c_L$ comes from a Poincar\'e inequality for one step, in which the gauge-fixing term is necessary (Lemma~\ref{lem:poincare-step}).
\item \emph{(Gauge fixing.)} $T^{\rm gf}(f)=T(f)$ for every $\mathcal G'$-invariant $f$ that vanishes outside Ba\l aban's set $\mathfrak M$ of regular configurations. Let $\chi=\chi^{\rm w}\chi'$ with a covariant cutoff $\chi'$, for instance $\chi'=1$ or $\chi'=\chi_\xi$. Then $T^{\rm gf}(\chi\rho)(V')$ does not depend on the choice of $U'$ in the minimal orbit of $V'$, and it is invariant under real gauge transformations of $V'$. If $\rho\ge0$, the part discarded by the cutoffs, $T(\one_{\mathfrak M\cap\mathfrak U_k(\tilde h)}\rho)-T^{\rm gf}(\chi\rho)$, is nonnegative.
\item \emph{(Lattice symmetry.)} $T$ commutes with the group $\mathfrak S_{k+1}$ generated by the translations of $\T'$ and by $B_4$ about the points of $\T'$. If $\rho\circ r=\rho$ for some $r\in\mathfrak S_{k+1}$, then $T^{\rm gf}(\chi\rho)(rV')=T^{\rm gf}(\chi\rho)(V')$ for the cutoffs $\chi$ of~(iii). Replacing $(W,U')$ by $(rW,rU')$ conjugates $\Delta_k$ and $C_k$ by an orthogonal map of the fields on $\mathfrak F$.
\end{enumerate}
The constants do not depend on $k$, $n$, $\eta$ or the coupling.
\end{maintheorem}

\paragraph{Theorem~\ref{thm:K} in words.} The step integrates over the configurations $V$ on $\T$ with prescribed averages $V'$ on $\T'$. Its base $W$ is the $k$-step average of the background field $U'=U_{k+1}(V')$; then $W$ has the averages $V'$, and $U'$ is also its $k$-step background field (Lemma~\ref{lem:composition}). On the support of $\chi^{\rm w}$, $V=e^{\ii Z}W$ with $|Z|\le w$ on every bond (Definition~\ref{def:chi-w}). The fluctuation variables are the values of $Z$ on $\mathfrak F$, and $\Delta_k$ is the quadratic part in $Z$ of the action with the gauge-fixing exponent (Lemma~\ref{lem:step-quadratic}). The density $\rho$ is $\exp(-g^{-2}S(U_k(V))+\mathfrak E_k(V))$ on the small-field set $\mathfrak U_k(\tilde h)$ \eqref{eq:rho-form}. So the step is an explicit prefactor, with $e^{-g^{-2}S(U')}$ and $(\det\Delta_k)^{-1/2}$, times a Gaussian integral with covariance $C_k=\Delta_k^{-1}$ of the cutoff and of the exponential of an interaction of third order, of Jacobian terms and of the old activities (i). The form $\Delta_k$ is coercive and $C_k$ decays exponentially, uniformly in $k$ (ii). The gauge fixing does not change the step on invariant integrands supported on regular configurations, and, for $\rho\ge0$, the cutoffs discard a nonnegative part (iii). The precise statements are Theorems~\ref{thm:step} and~\ref{thm:step-coercive}.

The lattice symmetry rests only on intrinsic objects: the Haar measure, the averaging, the coarse action and $\mathrm{gf}_W$. The gauge-field operators of Part~II at curved backgrounds are covariant under translations and permutations of the axes, but in general not under reflections (Remark~\ref{rem:symmetry-reflections}). Part~VII also shows that the curvature of background fields is bounded by $(C_*+\delta)$ times the curvature of the coarse configuration, with $C_*$ independent of $L$ (Proposition~\ref{prop:sharp-curvature}). This makes the small-field sets of consecutive steps compatible for $L^2\ge8(C_*+1)$, provided the couplings satisfy $g_{k+1}\le2g_k$ and are small enough that the fluctuation cutoff obeys $4w\le\frac12h_k$ (Remark~\ref{rem:parameters}, the condition (Q2)). Along the flow both hold at every step (Theorem~\ref{thm:N}(v)).

\paragraph{Localization (Parts VIII and IX).} The renormalization group needs the background dependence of the step as a sum of local pieces. Part~VIII obtains this with Dirichlet approximants and telescoping, using interior pointwise bounds with derivatives for the Dirichlet operators (Theorem~\ref{thm:ID}). The pieces are of \emph{radius type}: a piece depending on a ball of radius $m$ is bounded by $e^{-\gamma m}$. A polymer expansion needs \emph{volume type}, a small factor for every cube of the dependence region. Part~IX provides it, as Ba\l aban and Dimock do, by random walks with weakening parameters $s\in\mathbb D_\kappa$. It constructs them for the covariance, its square root, the coarse Hessian and the determinant in $L^2$ norms (Section~\ref{sec:nested-walks}). It proves interior estimates with $L^p$ data for the local propagators (Section~\ref{sec:interior-lp}). With them, regularity tiers give pointwise $C^{1,\beta}$ bounds for the weakened propagators $\bG(s)$, $\bH(s)$, $\bC(s)$ (Section~\ref{sec:tiers}). In four dimensions $L^2$ data lie at the endpoint $p=\frac d2$ of the interior estimates, and the tiers pass through block $L^4$ (Lemma~\ref{lem:endpoint-factors}). With these bounds Part~IX builds the weakened background map $A_*(s;B)$, analytic in $s$ and factorizing over polymers (Theorem~\ref{thm:weak-bg}).

\paragraph{The small-field step (Part X).} The activities are analytic functions on complex domains. In four dimensions the constants of these domains do not depend on the coupling, the step, the torus or the polymer (Remark~\ref{rem:domain-constants}), as in \cite[(1.11)--(1.16)]{BalYM4I}. Part~X first extends Part~IX, and the averaged backgrounds of Part~IV, to complex references (Section~\ref{sec:complex-references}). Then it follows Dimock's scheme \cite{DimockI}:
\begin{itemize}
\item the step quantities (interaction, gauge-fixing remainder, Jacobians, old activities composed with the background map, Hessian) become functions of the weakening parameters that factorize over polymers (Section~\ref{sec:step-quantities});
\item with the cutoffs on the fluctuation field and on the white noise, the fluctuation integral is evaluated by a Mayer and polymer expansion (Section~\ref{sec:cluster}). By parity in the coupling there is no term of order $g$; the pure sources are of order $g^2p_0^6$.
\end{itemize}
Here $p_0=p_0(g)$ is the white-noise cutoff, which grows more slowly than every power of $g^{-1}$, for instance as a power of $\log g^{-1}$ (Section~\ref{sec:flow}).

\begin{maintheorem}[The small-field step in polymer form; Theorem~\ref{thm:sf-step}, Propositions~\ref{prop:first-step} and~\ref{prop:vacuum}]\label{thm:L}
Let $G$ be compact, connected and semisimple, and let the H\"older exponent satisfy $\alpha\le\frac12$. Let the parameters be chosen as in Remark~\ref{rem:parameters}, and assume the conditions of Theorem~\ref{thm:sf-step} on the coupling, the thresholds and the old activities.
\begin{enumerate}[label=(\roman*)]
\item For a density $\rho$ of the form \eqref{eq:rho-form} at step $k$, the small-field step is, at real critical references with $V'$ small field,
\[
\begin{aligned}
T_{\rm sf}(\rho)(V')=\mathcal N_kN_\chi\exp\Big(&-\frac1{g^2}S(U')-\frac12\sum_b\ell_k(b;1)+\sum_X\mathfrak e(X,U')+\sum_c\log\mathfrak j_c\\
&+\mathfrak E(U')+\sum_YE^\#(Y,U')\Big),
\end{aligned}
\]
with explicit constants $\mathcal N_k$, $N_\chi$ and densities $\ell_k(b;1)$ at the trivial base. Every term depends on $U'$ only near its polymer, is invariant under real gauge transformations, is real at real references, and is analytic in the reference on complex domains. On the domains of Remark~\ref{rem:domain-constants} the one-loop activities are bounded by $Ce^{-(\kappa_2-1)|X|}$ and the Jacobian terms by $C$, with $C$ independent of the coupling; their quadratic parts are marginal in four dimensions. The new sources $E^\#=E^\#_{\rm p}+E^\#_{\rm m}$ consist of a pure part of order $g^2p_0^6$ and a mixed part bounded by the old activities $\mathfrak E$. If $\rho$ is invariant under some $r\in\mathfrak S_{k+1}$, so is the sum over polymers in the exponent; single activities are covariant at most under the translations of their grids.
\item The first step, from $\rho_0=\exp(-g_0^{-2}S-E_0)$, has the same form, with $\mathfrak E=0$ and no mixed sources.
\item \emph{(Vacuum energy.)} If the steps $k<K$ are given by (i) and (ii), with the hypotheses of (i) at every step, their constants split into explicit normalizations and a vacuum energy $\eps_k$ per site of $\T^{(k)}$, with $\eps_{k+1}=L^d\eps_k+\delta_k$, where $\delta_k$ does not depend on $E_0$, its pure part is at most $CL^dM^{-d}g_k^2p_0(g_k)^6$, and its mixed part is bounded by the old activities. There is exactly one counterterm $E_0$ with $\eps_K=0$. If $|\delta_{k+m}|\le D_kL^{dm/2}$ for $0\le m<K-k$, then $|\eps_k|\le2L^{-d}D_k$. Along the flow of Theorem~\ref{thm:N} the hypotheses of~(i) hold at every step, and the growth bound holds with $D_k=Cg_kp_0(g_k)$ (Theorem~\ref{thm:N}(vi)).
\end{enumerate}
\end{maintheorem}

In words, Theorem~\ref{thm:L} says the following. The small-field step maps a density of the form \eqref{eq:rho-form} to $e^{-g^{-2}S(U')}$ times explicit normalizations and densities and the exponential of a sum of activities on polymers. Each activity depends on the reference only near its polymer, is gauge invariant, and is analytic in the reference on the complex domains of Remark~\ref{rem:domain-constants}, whose constants do not depend on the coupling (i). The new activities are of three kinds. The one-loop activities and the Jacobian terms are bounded independently of the coupling, and their quadratic parts are marginal. The pure sources are of order $g^2p_0^6$; by parity in the coupling there is no term of order $g$ (Proposition~\ref{prop:pure-sources}). The mixed sources are bounded by the old activities. The hypotheses on the coupling, the thresholds and the old activities hold at every step of the flow (Theorem~\ref{thm:N}(v)). The small-field steps define a modified model: at every step the integration is restricted to small fields, and the rest is discarded. The discarded parts are nonnegative, so the partition function of the small-field model is at most that of the lattice theory with the same counterterm (Remark~\ref{rem:sf-reference}). Theorem~\ref{thm:L} is a statement about one step; it is not a flow theorem. The iteration needs the four-dimensional organization of the flow \cite{BalYM4I,BalYM4II}: generations of activities, each created by one step and bounded at its creation scale; the cancellation of the marginal quadratic parts of old generations by counterterms $-\beta_jS$, with $g_{j+1}^{-2}=g_j^{-2}-\beta_{j+1}$; and an iteration without contraction. Part~X carries this out: the generations, their symmetrization and their coefficient $\beta_j$ (Section~\ref{sec:generations}), the extraction of the marginal part of an old generation at later scales (Section~\ref{sec:extraction}), one step in generation form (Theorem~\ref{thm:gen-step}), and the flow over all scales at which the couplings stay in $(0,\gamma]$ (Theorem~\ref{thm:sf-flow}). The flow supplies, at every step, the bounds $g_{k+1}\le2g_k$ and the bounds on the old activities that Theorem~\ref{thm:L} assumes, and the growth bound on the $\delta_k$ in (iii). It is summarized in Theorem~\ref{thm:N} below: Theorem~\ref{thm:L} is the step, Theorem~\ref{thm:N} the flow. Controlling the discarded parts is the large-field stage, the analogue of \cite{DimockII,DimockIII}, which is not treated here. Parts~VII--X write out the small-field part of Ba\l aban's four-dimensional analysis \cite{BalYM4I,BalYM4II}, for one step and for its iteration over all scales on a fixed torus, as long as the couplings stay in $(0,\gamma]$, in Dimock's scheme, on the operator theory of Parts~I--IX.

\emph{Boundary conditions.} The gauge-field Dirichlet operators contain an exterior divergence term, so that all components of the gauge field satisfy Dirichlet conditions (Section~\ref{sec:reg-gauge}). With these conditions the local commutators of the weakened walks are bounded uniformly in $\eta$ on proper domains (Remark~\ref{rem:L3-domains}). On the torus the exterior term vanishes, so the torus operators, and with them Theorems~\ref{thm:F}, \ref{thm:H}--\ref{thm:J} and the torus statements of Theorem~\ref{thm:G}, do not involve it.

\paragraph{The one-loop coefficient (Part XI).} In four dimensions the coupling is marginal, and its flow is driven by the coefficient $\beta$. Part~XI computes the one-loop part of this coefficient, summed over the steps. Let $\Gamma^{(1)}_k(V)$ be the one-loop term of the composite $k$-step map, obtained by Laplace's method from the integral over the fine configurations with averages $V$ near the minimal orbit (Proposition~\ref{prop:ol-laplace}), and $\Gamma^{\rm step}_j$ the one-loop term of the step from scale $j$ to scale $j+1$ (Section~\ref{sec:ol-composite}). For $j\ge1$ it is given by Theorem~\ref{thm:K}(i). For $j=0$ it is given by Lemma~\ref{lem:first-step-formula}, the formula of Theorem~\ref{thm:K}(i) for the step from scale $0$ with the fluctuation operator of Proposition~\ref{prop:first-step}(b), the \emph{formula for the first step}. For a gauge-invariant function $\Phi$ of the coarse configuration, $\Pi_\Phi(p)$ is its vacuum polarization, the second variation at $V=1$ in momentum space, and $T_\Phi$ its marginal tensor, the part of $\Pi_\Phi(p)$ of second order in $p$ (Definition~\ref{def:marginal-tensor}). Theorem~\ref{thm:M} is stated before Theorem~\ref{thm:N}, and its part~(iv) uses the flow of Theorem~\ref{thm:N}; the generations and the couplings of the flow are introduced in the paragraph on Part~X before Theorem~\ref{thm:N}.

\begin{maintheorem}[The one-loop coefficient; Lemma~\ref{lem:ol-one-step}, Proposition~\ref{prop:ol-composition}, Lemma~\ref{lem:first-step-formula}, Theorem~\ref{thm:ol-chart}, Lemmas~\ref{lem:ol-symmetry}, \ref{lem:ol-torus-limit} and~\ref{lem:ol-limits}, Theorems~\ref{thm:ol-log}, \ref{thm:ol-cumulative} and~\ref{thm:ol-uv}, Proposition~\ref{prop:flow-F}]\label{thm:M}
Let $G$ be compact, connected and simple, with the trace of a representation $\rho$ in the action, $\operatorname{tr}_\rho(T^aT^b)=T_\rho\delta^{ab}$, and $C_A$ the Casimir of the adjoint representation.
\begin{enumerate}[label=(\roman*)]
\item \emph{(Composition and chart.)} For every small-field $V$ and $k\ge2$, up to a constant independent of $V$, $\Gamma^{(1)}_k(V)=\Gamma^{(1)}_{k-1}(U_k(V)^{(k-1)})+\Gamma^{\rm step}_{k-1}(U_k(V))$. It follows that $\Gamma^{(1)}_k(V)=\sum_{j<k}\Gamma^{\rm step}_j(U_k(V))$ for $k\ge1$, up to constants independent of $V$. Moreover, $\Gamma^{(1)}_k(V)$ is given in the Landau chart at the background field $U_k(V)$ by the determinants of $\bL_*$, of $\bQ_k\bG_*\bQ_k^T$ and of a ghost operator.
\item \emph{(Symmetry.)} $\Gamma^{(1)}_k$ and the functions $V'\mapsto\Gamma^{\rm step}_j(U_{j+1}(V'))$, $j\ge0$, are exactly invariant under the translations and under $B_4$. As the side of the torus tends to infinity along the admissible sides (Lemmas~\ref{lem:ol-torus-limit} and~\ref{lem:ol-limits}), their second variations at $1$ converge, with exponential decay uniform in the side. In infinite volume, $\Pi(Rp)=R\,\Pi(p)R^T$ for $R\in B_4$, $\Pi(p)=O(|p|^2)$, and the marginal tensor is $\beta(\delta_{\mu\nu}|p|^2-p_\mu p_\nu)$ with a single coefficient $\beta$; in particular there is no term linear in $p$. These functions are invariant under single reflections, so no topological density can occur in their local expansions.
\item \emph{(The cumulative coefficient.)} The marginal tensors $T^{(1)}_{j+1}$ of the steps satisfy, uniformly in $k$,
\[
\sum_{j<k}T^{(1)}_{j+1}=\big(k\beta_\infty+O(1)\big)\big(\delta_{\mu\nu}|p|^2-p_\mu p_\nu\big),\qquad\beta_\infty=\frac{11C_A}{24\pi^2T_\rho}\log L .
\]
\item \emph{(The ultraviolet limit.)} Let $p_0(g)=(\log g^{-1})^q$, let the parameters and the torus be as in Theorem~\ref{thm:N}, and let $\gamma$ be as in Theorem~\ref{thm:ol-uv}. Then there are $g_{\rm uv}>0$ and $C$, independent of $K$ and of the torus, such that for $0<g\le g_{\rm uv}$ and every $K$ the small-field flow of Theorem~\ref{thm:N} has a bare coupling $g_0$, depending on $K$, $g$ and the torus, with $g_k\in(0,\gamma]$ for $k\le K$, $g_K=g$, and
\[
\Big|\frac1{g_k^2}-\frac1{g^2}-(K-k)\beta_\infty\Big|\le C\big(\log(K-k+2)\big)^{6q+1}\qquad\text{for every }k\le K.
\]
The properties (F1)--(F3) of the flow that this uses are proved in Proposition~\ref{prop:flow-F}. Property (F2), the comparison of the one-loop parts of the coefficients $\beta_{j+1}$ of the flow with the coefficients of~(iii), rests on the torus independence of the one-loop pieces (Proposition~\ref{prop:one-loop-local}, Lemma~\ref{lem:OL}). Here $g$ is the coupling of the last scale; in Parts~VII and~X, $g=g_0$ denotes the bare coupling.
\end{enumerate}
\end{maintheorem}

\paragraph{Theorem~\ref{thm:M} in words.} The one-loop term of $k$ steps is the sum of the one-loop terms of the single steps, up to constants, and it is given in the Landau chart by three determinants (i). These functions of the coarse configuration are invariant under the lattice symmetries. Hence the infinite-volume limits of their second variations at $1$ have no term linear in $p$, and their marginal tensors are multiples of $\delta_{\mu\nu}|p|^2-p_\mu p_\nu$ (ii). Summed over $k$ steps, the coefficient is $k\beta_\infty+O(1)$ (iii). With the flow of Theorem~\ref{thm:N} this gives~(iv): for a small final coupling $g$ and any number $K$ of steps there is a bare coupling that leads to $g_K=g$ through couplings in $(0,\gamma]$, and $g_k^{-2}-g^{-2}$ is $(K-k)\beta_\infty$ up to a power of $\log(K-k+2)$. Since $\beta_\infty>0$, at $k=0$ this says that $g_0^{-2}$ grows like $K\beta_\infty$: at a fixed $g$ the bare coupling tends to $0$ as the number of steps grows. Part~(iv) concerns the couplings of the small-field flow (Remark~\ref{rem:ol-remains}(a)).

Part~(ii) rests on the symmetry of Ba\l aban's averaging, the invariance of the Haar measure and the covariance of the minimal orbits. It does not use the covariance of the gauge-field operators of Part~II, which are covariant under reflections only at flat backgrounds. For the step functions it has two proofs: one through the invariance of $\Gamma^{(1)}_k$, which uses no covariance property of the gauge fixing, and a direct one, which uses only the covariance of the quadratic part of the gauge fixing, that of the averaged contour variables. In particular there is no term linear in $p$; such a term would produce relevant terms of order $1$ in every step (Remark~\ref{rem:marginal-orders}, Remark~\ref{rem:ol-remains}(b)). The theorems of Part~XI have no hypothesis on a linear term. Parts~(iii) and~(iv) are our version, for the couplings of the small-field flow of Part~X, of Ba\l aban's Theorem~2 \cite{BalYM4I}, which he states for $SU(2)$. Ba\l aban's couplings come from his own generations, a different construction (Remark~\ref{rem:sf-flow}(a) and the statement of Proposition~\ref{prop:beta-identification}). Part~(iii) is the growth $k\beta_\infty+O(1)$, in particular the positivity for large $k$, of the one-loop coefficient summed over the steps, and (iv) derives from it the behaviour of the couplings at all scales, with the flow of Part~X (Theorem~\ref{thm:sf-flow}, Proposition~\ref{prop:flow-F}) in place of his Theorem~3. In particular (iv) gives the analogue of his bounds (0.31) for these couplings at the scales $k$ with $K-k\ge m_1$, for some $m_1$ independent of $K$, $g$ and the torus; for the last $m_1$ steps it gives the upper bound, and the lower bound up to an additive constant (the paragraph after Theorem~\ref{thm:ol-uv}). The main input of (iii) is the coefficient $\beta_{\rm F}=k\beta_\infty+O(1)$ of the vacuum polarization on the fine lattice (Theorem~\ref{thm:ol-log}). The infinite-volume objects of Part~XI are constructed from the tori: the operators at $U=1$ by folding (Section~\ref{subsec:ol-infinite-volume}), and the one-loop pieces of the steps through their torus independence (Proposition~\ref{prop:one-loop-local}). Remark~\ref{rem:ol-remains} lists what remains.

\paragraph{The small-field flow (Part X).} The flow iterates the step of Theorem~\ref{thm:L} in Ba\l aban's generation form \cite[(0.22)--(0.30)]{BalYM4I}, on a torus whose side at the last scale $K$ is $n_K\ge n_{\min}$ (Lemma~\ref{lem:torus}). A \emph{generation} of scale $j$ is a family $E=\{E(X,\cdot)\}$ of normalized, gauge-invariant analytic activities on the polymers $X$ of the grid of scale $j$. They are real at real configurations and covariant under the group $\mathfrak S^{\rm grid}_j$ of the symmetries of that grid, and their sum at the background fields of scale $j$ is invariant under $\mathfrak S_j$ on the small-field set (Definition~\ref{def:generation}). A generation is bounded only at its own scale, on domains $\mathcal U^\flat(X)$ whose constants do not depend on the coupling, the step, the torus or the polymer. Its coefficient $\beta_E$ is read off from its period density $\bar T_E$, the second variation of its activities at the linear fields of constant curvature, per unit volume, as a multiple of that of $S$ (Definition~\ref{def:period-density}, Proposition~\ref{prop:period-iso}). The new terms of a step, i.e.\ the one-loop activities, the Jacobian terms and the new sources of Theorem~\ref{thm:L}(i), reblocked, form a raw generation (Lemma~\ref{lem:raw-new}); the new generation is its average over $\mathfrak S^{\rm grid}_{k+1}$ (Definition~\ref{def:symmetrization}). The old generations enter the next step through the old activities of Definition~\ref{def:old-activities}: in every grid cube $\Box$ the counterterm $-\beta_jS_\Box$ is grouped with the small polymers of the generation $j$ in $\Box$. We write $(E-\beta S)(U)=\sum_XE(X,U)-\beta S(U)$, and $h_k=p(g_k)g_k$, $p_{0,k}=p_0(g_k)$ with the small-field parameters of Section~\ref{subsec:flow-parameters}.

\begin{maintheorem}[The four-dimensional small-field flow; Theorem~\ref{thm:sf-flow}, Propositions~\ref{prop:period-iso}, \ref{prop:extraction} and~\ref{prop:flow-continuity}, Corollaries~\ref{cor:bare-continuity} and~\ref{cor:sf-stability}]\label{thm:N}
Let $G$ be compact, connected and simple, let the parameters be chosen as in Remark~\ref{rem:parameters}, with the bound $\gamma$ on the couplings subject to ($\Gamma$1)--($\Gamma$3) of Remark~\ref{rem:flow-constants}, and let $n_K\ge n_{\min}$. Let $g_0\in(0,\gamma]$, $E_0\in\R$ and $\rho_0=\exp(-g_0^{-2}S-E_0)$. For $k\ge0$, as long as $g_k\le\gamma$, let $E^{(k+1)}$ be the new generation of the step from scale $k$ to scale $k+1$ applied to $\rho_k$, let $\beta_{k+1}=\beta_{E^{(k+1)}}$ and
\[
\frac1{g_{k+1}^2}=\frac1{g_k^2}-\beta_{k+1},
\]
and, if $g_{k+1}\le\gamma$, let $\rho_{k+1}=\one_{\mathfrak U_{k+1}(h_{k+1})}T_{\rm sf}(\rho_k)$. There are constants $\mathsf K$, $C_\beta$, $C$ and $C'$, which depend only on the parameters fixed before $\gamma$ ($C'$ also on the function $p_0$), and not on $k$, $K$, the torus or the couplings, such that (i)--(v) hold for every $1\le k\le K$ with $g_1,\dots,g_k\le\gamma$.
\begin{enumerate}[label=(\roman*)]
\item \emph{(Generation form.)} For $V\in\mathfrak U_k(h_k)$, with $U_k=U_k(V)$,
\[
\rho_k(V)=\mathcal Z_k\exp\Big(-\frac1{g_k^2}S(U_k)-\eps_k|\T^{(k)}|+\sum_{j=1}^k\big(E^{(j)}-\beta_jS\big)(U_k)\Big),
\]
and $\rho_k(V)=0$ for the other configurations $V$. The density $\rho_k$ is invariant under real gauge transformations and under $\mathfrak S_k$.
\item \emph{(Creation scale.)} Each $E^{(j)}$, $j\le k$, is a generation of scale $j$, and $\sup_{\mathcal U^\flat(X)}|E^{(j)}(X,\cdot)|\le\mathsf Ke^{-\kappa_nd_j(X)}$ for every polymer $X$ of scale $j$.
\item \emph{(Coefficients.)} $\bar T_{E^{(j)}}(F)=\beta_j\bar T_S(F)=\frac12\beta_j\sum_{\mu<\nu}\operatorname{tr}\big(F_{\mu\nu}^2\big)$ for all curvature data $F$: the period density has no topological and no anisotropic part. The number $\beta_j$ is real, and $|\beta_j|\le C_\beta\mathsf K$.
\item \emph{(Later scales.)} Let $j\le k$, $m=k-j$, and let $\Box$ be a grid cube of scale $k$. The sum of $E^{(j)}(X,U)$ over the small polymers $X$ of scale $j$ based in $\Box$ (Definition~\ref{def:polymer-classes}) differs from $\beta_jS_\Box(U)$ by at most $C\mathsf KL^{-\alpha m}$ if $m\ge m_0$, and by at most $C\mathsf K$ if $m<m_0$, for $U\in\mathcal U^\flat(\Box^+)$ of scale $k$. The sum over the other polymers of scale $j$ whose closure at scale $k$ is $Y$ is at most $C\mathsf KL^{-m}e^{-\kappa_nd_k(Y)}$ on $\mathcal U^\flat(Y)$. In particular $|\sum_{j=1}^k(E^{(j)}-\beta_jS)(U_k(V))|\le C\mathsf K|\T^{(k)}|$ for $V\in\mathfrak U_k(h_k)$.
\item \emph{(Couplings.)} For $1\le j\le k$, $|g_j^{-2}-g_{j-1}^{-2}|=|\beta_j|\le C_\beta\mathsf K$ and $\frac34g_{j-1}\le g_j\le2g_{j-1}$, and the step from scale $j-1$ to scale $j$ satisfies the hypotheses of Theorem~\ref{thm:L}(i) (of Theorem~\ref{thm:L}(ii) if $j=1$).
\item \emph{(Vacuum and stability.)} Let $g_1,\dots,g_K\le\gamma$. The constants $\mathcal Z_k$ are explicit products of normalizations. The vacuum energies satisfy $\eps_{k+1}=L^4\eps_k+\delta_k$ for $k<K$, where $\delta_k$ does not depend on $E_0$ and $|\delta_k|\le C(g_k^2p_{0,k}^6+g_kp_{0,k})$. There is exactly one $E_0$ with $\eps_K=0$, and for it $|\eps_k|\le C'g_kp_{0,k}$ for $k\le K$. The partition function of the small-field model is $\int\rho_K=\mathcal Z_Ke^{-\eps_K|\T^{(K)}|}I_K$, where $I_K$ does not depend on $E_0$ and
\[
e^{-c|\T^{(K)}|}\le I_K\le e^{c|\T^{(K)}|},
\]
with $c$ depending on $g_K$ and on the parameters, but not on $K$, the torus or $g_0,\dots,g_{K-1}$.
\item \emph{(Continuity.)} Let $0\le j<K$. For every sequence $(g_0,\dots,g_j)\in(0,\gamma]^{j+1}$ with $g_{i+1}\le2g_i$ for $i<j$, the generation $E^{(j+1)}$ and the coefficient $\beta_{j+1}$ are defined by the steps above, with the couplings taken from the sequence and not from the recursion, as in Proposition~\ref{prop:flow-continuity}. They depend only on $(g_0,\dots,g_j)$, and continuously, at every point of the domains of the activities. For the couplings of the recursion they are those of the preamble. The set of the $g_0\in(0,\gamma]$ with $g_1,\dots,g_K\le\gamma$ is closed in $(0,\gamma]$ and contains $(0,(\gamma^{-2}+KC_\beta\mathsf K)^{-1/2}]$, and $g_0\mapsto(g_1,\dots,g_K)$ is continuous on it.
\end{enumerate}
\end{maintheorem}

\paragraph{Theorem~\ref{thm:N} in words.} As long as the couplings stay in $(0,\gamma]$, the density of every scale keeps Ba\l aban's generation form (i). Each step creates one generation, a family of activities that is bounded once, at its creation scale, on domains whose constants do not depend on the coupling (ii). Its coefficient $\beta_j$ is read off from its period density, which is exactly isotropic (iii). At a later scale, $m\ge m_0$ scales after its creation, the small polymers of a generation in a grid cube add up to $\beta_jS_\Box$ up to $C\mathsf KL^{-\alpha m}$, and its other polymers carry a factor $L^{-m}$; so the counterterm $-\beta_jS$ cancels its marginal part, and the sum over all generations is bounded per unit volume, uniformly in the scale (iv). The iteration uses no contraction (Remark~\ref{rem:flow-status}). The couplings change by bounded amounts per step, and each step satisfies the hypotheses of Theorem~\ref{thm:L} (v). Item~(vi) writes, for the counterterm $E_0$ with $\eps_K=0$, the partition function of the small-field model as the explicit product $\mathcal Z_K$ of normalizations times a number between $e^{-c|\T^{(K)}|}$ and $e^{c|\T^{(K)}|}$, with $c$ depending on $g_K$ but not on the number of steps or the torus. Item~(vii) gives the continuity in the couplings that the exit argument in the proof of Theorem~\ref{thm:ol-uv} uses. Section~\ref{g:result} restates the theorem in more detail.

\paragraph{Concordance with Section~\ref{sec:flow4}.} Items (i)--(vi) of Theorem~\ref{thm:N} are items (a)--(f) of Theorem~\ref{thm:sf-flow}, in this order. The partition function in~(vi) is Corollary~\ref{cor:sf-stability}, and item~(vii) is Proposition~\ref{prop:flow-continuity} with Corollary~\ref{cor:bare-continuity}. Item~(iii) is also Proposition~\ref{prop:period-iso}, and item~(iv) is Proposition~\ref{prop:extraction}, Lemma~\ref{lem:large-later}(a) and Remark~\ref{rem:old-later}. The constant $C$ of Theorem~\ref{thm:N} stands for $C_{\rm ex}$, $C_{\rm ne}$ and $C_{\rm lat}$ in~(iv) and for $C_{\rm vac}$ in~(vi), and $C'$ is $C'_{\rm vac}$ (Theorem~\ref{thm:sf-flow}, Lemma~\ref{lem:large-later}).

Theorem~\ref{thm:N} is the analogue, for the small-field model on the torus, of Ba\l aban's Theorem~3 \cite{BalYM4I}, with his hypothesis $0<g_k\le\gamma$ for $k\le K$; as there, no asymptotic behaviour of the couplings is assumed (Remark~\ref{rem:sf-flow}(a)). Unlike there, the dependence on the couplings, which is a dependence on the whole sequence $(g_0,\dots,g_{j-1})$, is only shown to be continuous, whereas his is smooth, and the small-field domain shrinks with the coupling, as in Dimock's scheme (Remark~\ref{rem:sf-flow}(a)). Items (i), (ii) and (iv) correspond to his (0.22)--(0.24), (0.25) and (0.28)--(0.30); the later-scale bounds (0.29) are asserted by his Theorem~1, and they are not among the inductive assumptions listed in his Theorem~3. The proof is an induction over the scales with the step in generation form (Theorem~\ref{thm:gen-step}). Its two four-dimensional ingredients are the symmetrization of every new generation over the symmetries of its grid, which makes the period density exactly isotropic (Proposition~\ref{prop:period-iso}), and the extraction of the marginal part of an old generation at every later scale. For the latter, a configuration of a later scale is written, in a suitable gauge on each small polymer of the generation, as a linear field of constant curvature plus a smaller remainder (Lemma~\ref{lem:later-normal-form}, Proposition~\ref{prop:extraction}). The old activities are then bounded uniformly in $k$ (Proposition~\ref{prop:old-activities}), and the iteration needs no contraction. The coefficients $\beta_j$ are defined on the given torus. Their one-loop parts are compared with the coefficients of Theorem~\ref{thm:M}(iii) in (F2) of Proposition~\ref{prop:flow-F} (Theorem~\ref{thm:M}(iv)). Ba\l aban's definition \cite[(1.20)--(1.22)]{BalYM4I} assigns a coefficient to the family of the generations $E^{(j)}$ of this note built with fixed couplings on tori of growing side. Proposition~\ref{prop:beta-identification} shows that this coefficient is defined and that, at fixed couplings, $\beta_j$ converges to it as the torus grows. This rests on the generation locality of Section~\ref{sec:gen-locality}; no rate is proved. Theorem~\ref{thm:M} does not use it.

\paragraph{Reading guide.} Section~\ref{g:read} of the guide gives four reading paths, Figure~\ref{fig:G-mainline} the main line, and the introduction of each Part its imports and exports. The main line is as follows. For the small-field step it is Parts~I--VII, Section~\ref{sec:dirichlet-averaged} and the interior estimates of Section~\ref{sec:interior}, Section~\ref{sec:localization} up to Proposition~\ref{prop:local-fluct}, Section~\ref{sec:det-local} except Proposition~\ref{prop:det-local} and Corollary~\ref{cor:one-loop}, Part~IX and Part~X. The radius-type results of Part~VIII (Theorem~\ref{thm:local-covariance}, Corollary~\ref{cor:large-changes}, Proposition~\ref{prop:det-local}, Corollary~\ref{cor:one-loop}, and Section~\ref{sec:interaction-local}) are superseded on the main line by Parts~IX and~X. Each of them is marked as superseded where it appears, in the status paragraphs at the head of Sections~\ref{sec:localization}--\ref{sec:interaction-local}. The flow (Theorem~\ref{thm:N}) is Sections~\ref{sec:flow}--\ref{sec:flow4} of Part~X. Sections~\ref{sec:ol-composite}--\ref{sec:ol-feynman} of Part~XI use Parts~I--VII and, from Part~X, mainly Propositions~\ref{prop:jacobians} and~\ref{prop:first-step} and Lemma~\ref{lem:first-step-formula}; Section~\ref{sec:ol-uv}, which combines Part~XI with the flow, uses Part~X throughout, and the torus independence of the one-loop pieces (Section~\ref{subsec:torus-independence}). Section~\ref{sec:gen-locality} extends it to whole steps at fixed old data, at flat configurations, and proves that the generations converge as the torus grows, at fixed couplings; it uses Parts~IV--X, Lemma~\ref{lem:ol-folding}, and Sections~\ref{subsec:ol-flow-coeff} and~\ref{subsec:ol-period-vacpol}.

\subsection{Methods}\label{ss:1.4}

Table~\ref{tab:tools} lists the main tools of the note, Part by Part, with the place where each is introduced and what it gives. Section~\ref{g:hard} of the guide names the places where the work is hardest.

\begin{table}[tbp]
\centering\footnotesize
\renewcommand{\arraystretch}{1.12}
\begin{tabular}{|p{0.06\textwidth}|p{0.30\textwidth}|p{0.22\textwidth}|p{0.31\textwidth}|}
\hline
Part & Tool & Introduced in & What it gives\tabularnewline
\hline
I & \raggedright block Poincar\'e inequality; for gauge fields the near separation of gradient and co-closed parts by $\bQ$ & \raggedright Section~\ref{sec:poincare}, Lemma~\ref{lem:Q-divergence} & \raggedright coercivity with an averaging term (Theorem~\ref{thm:A})\tabularnewline
I & \raggedright the Combes--Thomas method & \raggedright Section~\ref{sec:decay} & \raggedright exponential decay of the inverses, on the fine and on the coarse lattice (Theorems~\ref{thm:B}--\ref{thm:E})\tabularnewline
I & \raggedright a lattice Schauder estimate & \raggedright Theorem~\ref{thm:local-regularity}, Appendix~\ref{app:green} & \raggedright pointwise bounds on values, derivatives and H\"older quotients from $L^2$ bounds\tabularnewline
I & \raggedright one Fourier lower bound & \raggedright Corollary~\ref{cor:QKQ-lower} & \raggedright the lower bound on $QGQ^*$, uniformly in $k$\tabularnewline
II & \raggedright perturbation in the background & \raggedright Section~\ref{sec:perturbation} & \raggedright the operators at backgrounds close to $1$ in a global gauge (Theorem~\ref{thm:F})\tabularnewline
III & \raggedright IMS localization at scale $M$; Dirichlet operators; comparison of domains & \raggedright Theorems~\ref{thm:reg-coercive} and~\ref{thm:reg-locality} & \raggedright coercivity at locally regular backgrounds, and locality in the background (Theorem~\ref{thm:G})\tabularnewline
IV & \raggedright the averaged axial gauge and a recursive correspondence map & \raggedright Theorems~\ref{thm:correspondence} and~\ref{thm:chart} & \raggedright the axial-to-Landau chart (Theorem~\ref{thm:H})\tabularnewline
V & \raggedright a contraction around a regular reference & \raggedright Proposition~\ref{prop:contraction} & \raggedright the variational problem relative to a reference (Theorem~\ref{thm:I})\tabularnewline
VI & \raggedright a change of the reference near a point; patching of local gauges & \raggedright Lemmas~\ref{lem:improvement} and~\ref{lem:large-cubes} & \raggedright the induction over $k$ on the torus (Theorem~\ref{thm:J})\tabularnewline
VII & \raggedright exponential gauge fixing without a Faddeev--Popov determinant; a Poincar\'e inequality for one step & \raggedright Lemmas~\ref{lem:exp-gauge} and~\ref{lem:poincare-step} & \raggedright the step as a Gaussian integral, and its coercivity (Theorem~\ref{thm:K})\tabularnewline
VIII & \raggedright Dirichlet approximants and telescoping; absorption inside a cube & \raggedright Theorem~\ref{thm:ID}, Section~\ref{sec:interior}, Lemma~\ref{lem:telescoping} & \raggedright interior bounds; localization of radius type (the telescoping is off the main line)\tabularnewline
IX, X & \raggedright random walks with weakening parameters; the localization lemma & \raggedright Theorem~\ref{thm:weak-inverse}, Lemma~\ref{lem:localization} & \raggedright localization of volume type\tabularnewline
IX & \raggedright interior $L^p$ estimates and regularity tiers & \raggedright Sections~\ref{sec:interior-lp} and~\ref{sec:tiers} & \raggedright pointwise $C^{1,\beta}$ bounds for the weakened propagators\tabularnewline
X & \raggedright complex references & \raggedright Proposition~\ref{prop:complex-ext} & \raggedright analyticity of the activities in the background\tabularnewline
X & \raggedright Mayer and polymer expansion with two cutoffs; parity in the coupling & \raggedright Theorem~\ref{thm:mayer}, Proposition~\ref{prop:pure-sources} & \raggedright the step in polymer form, with pure sources of order $g^2p_0^6$ (Theorem~\ref{thm:L})\tabularnewline
X & \raggedright symmetrization of a generation; a normal form at later scales & \raggedright Lemma~\ref{lem:symmetrization}, Lemma~\ref{lem:later-normal-form}, Proposition~\ref{prop:extraction} & \raggedright one coefficient $\beta_j$ per generation; old activities bounded uniformly in the scale, without contraction (Theorem~\ref{thm:N}, Proposition~\ref{prop:old-activities})\tabularnewline
XI & \raggedright Laplace's method; the lattice symmetry of the averaging & \raggedright Propositions~\ref{prop:ol-laplace} and~\ref{prop:ol-composition}, Lemma~\ref{lem:ol-symmetry} & \raggedright the one-loop term as a sum over the steps, with no term linear in $p$ (Theorem~\ref{thm:M}(i),~(ii))\tabularnewline
XI & \raggedright folding of the torus operators; torus independence of the one-loop pieces & \raggedright Section~\ref{subsec:ol-infinite-volume}, Proposition~\ref{prop:one-loop-local} (Part~IX) & \raggedright the operators in infinite volume; the comparison (F2) of the flow with Part~XI (Proposition~\ref{prop:flow-F})\tabularnewline
XI & \raggedright reduction to the lattice Feynman gauge & \raggedright Proposition~\ref{prop:ol-reduction}, Theorem~\ref{thm:ol-log} & \raggedright the coefficient $\beta_\infty$ (Theorem~\ref{thm:M}(iii))\tabularnewline
XI & \raggedright an exit argument & \raggedright Theorem~\ref{thm:ol-uv}, Corollary~\ref{cor:bare-continuity} & \raggedright the bare coupling of Theorem~\ref{thm:M}(iv)\tabularnewline
\hline
\end{tabular}
\caption{The main tools of the note.}\label{tab:tools}
\end{table}

Part~I uses four tools. Each is introduced in its own section, before it is used.

\paragraph{1. Coercivity (Section~\ref{sec:poincare} and Section~\ref{sec:gauge-coercivity}).} A Poincar\'e inequality on unit blocks shows that the averaging term controls exactly those modes that the Laplacian does not control. This gives the bounds of Theorem~\ref{thm:A} for scalars. For gauge fields the situation is subtler, because $\|dA\|^2+a\|\bQ A\|^2$ is invariant under gauge transformations $A\mapsto A+d\lambda$ with $Q_0\lambda=0$. The gauge-fixing term $\|Rd^*A\|^2$ removes this degeneracy. We prove the coercivity in position space. We split $A$ into a gradient part and a co-closed part, and show that the averaging operator nearly separates them: $\dl^*\bQ A$ differs from an average of $d^*A$ by the unit-lattice divergence of a quantity bounded by $\|\partial A\|$ (Lemma~\ref{lem:Q-divergence}). Ba\l aban's original proof \cite{BalPropI} is a computation in momentum space.

\paragraph{2. The Combes--Thomas method (Section~\ref{sec:decay}).} If an operator is coercive and conjugation by $e^{\alpha\rho}$, for Lipschitz $\rho$, perturbs it by $O(\alpha)$, then its inverse decays exponentially \cite{CT}. This works for the fine-lattice operators $G_\sigma$ and $\bG$, which contain the large term $-\Lap$ of norm $\sim\eta^{-2}$, because the perturbation is controlled by the quadratic form itself. It also works for the coarse operators $(QGQ^*)^{-1}$. Ba\l aban and \cite{DST1} use Fourier analysis or random-walk expansions instead. We use neither: on a single torus at $U=1$, Combes--Thomas suffices, and it also survives perturbation in~$U$.

\paragraph{3. Local regularity (Section~\ref{sec:regularity}).} Exponential decay is first obtained in $L^2$ between unit blocks. To pass to pointwise bounds on values and derivatives, uniformly in $\eta$, we prove a discrete interior regularity estimate. If $-\Lap u=F$ holds on a cube of side $4$ (with $u$ defined on a slightly larger cube), then $u$, $\partial u$ and the H\"older quotients of $\partial u$ on the concentric cube of side $2$ are bounded by $|F|$ and the $L^2$ norm of $u$. The proof uses a cutoff and the decay of the Green function of the fine lattice $\eta\Z^d$ and its derivatives. The latter is proved in Appendix~\ref{app:green}.

\paragraph{4. One Fourier lower bound (Section~\ref{sec:fourier}).} Averaged Green functions $QGQ^*$ must be bounded \emph{below} uniformly in $k$. We deduce this from the operator inequality $G\ge(-\Lap+2a)^{-1}$ and an elementary computation in momentum space. The same computation gives one more estimate needed in the gauge-field coercivity proof.

\subsection{What is not done here}\label{ss:1.5}

The variational theorem is proved on a single torus with the small-field condition everywhere (Theorem~\ref{thm:J}). The renormalization group is carried out for one small-field step (Theorems~\ref{thm:K} and~\ref{thm:L}) and iterated over all scales in the small-field model, on a fixed torus, as long as the couplings stay in $(0,\gamma]$ (Theorem~\ref{thm:N}). Two omissions matter most.
\begin{itemize}
\item \emph{The large-field stage.} At every step the small-field model discards the large fields, and what it discards is nonnegative (Remark~\ref{rem:sf-reference}(c)). Controlling the discarded parts, and so removing the modification, is the large-field stage, the analogue of \cite{DimockII,DimockIII}; it is not treated (Remark~\ref{rem:sf-reference}(d)). Hence nothing is proved about the partition function or the correlation functions of the lattice theory itself.
\item \emph{The continuum limit.} No limit of the measure or of correlation functions as the lattice spacing tends to $0$ is considered. Theorem~\ref{thm:M}(iv) is a statement about the couplings of the small-field flow, uniform in the number of steps at a fixed last torus (Theorem~\ref{thm:ol-uv}). With $K=N$ steps on a torus of physical side $\ell$ the bounds of Theorems~\ref{thm:N} and~\ref{thm:M}(iv) are uniform in the spacing $L^{-N}$ (Remark~\ref{rem:sf-flow}(d)); they are not statements about a limit of the measure.
\end{itemize}
In addition, we do not treat:
\begin{itemize}
\item a rate in the comparison of the $\beta_j$ with the coefficient that Ba\l aban's (1.20)--(1.22) assigns to the generations of this note: Proposition~\ref{prop:beta-identification} proves the convergence as the torus grows, at fixed couplings, without a rate and without uniformity in the couplings (Remark~\ref{rem:gl-rate}). Theorem~\ref{thm:M} does not use it;
\item the exact lower half of the analogue of Ba\l aban's (0.31) for the last $m_1$ steps, which Theorem~\ref{thm:M}(iv) gives up to an additive constant (Remark~\ref{rem:ol-remains}(c));
\item the nested domains $\Omega_j$ of \cite{BalVar}, which belong to the large-field stage;
\item gauges on cubes of size up to $c\,\eps_1^{-1}$ with bounds linear in the size \cite[(9), (10)]{BalVar}; we obtain polynomial bounds (Remark~\ref{rem:large-cubes});
\item gauge groups with abelian factors in Part~X: the second-order bound for normalized activities (Lemma~\ref{lem:quadratic-bound}) uses semisimplicity;
\item dimensions other than $4$.
\end{itemize}
Section~\ref{subsec:checklist} lists the input of the variational theorem and where it is provided.

\paragraph{Relation to \cite{DST1}.} The paper \cite{DST1} gives a self-contained account of exponential decay for the \emph{scalar} operators of Ba\l aban's approach on general domains with Neumann boundary conditions, using the method of images and analytic continuation of Fourier transforms. Here the domain is a torus, the methods are different, and the emphasis is on gauge fields. The scalar part (Sections~\ref{sec:scalar}--\ref{sec:scalar-avg}) is included because the gauge-field operators are built from scalar ones: gauge transformations are scalar fields, and the projection $R$ is defined in terms of $G_0$.

\paragraph{Status.} The note was written with the assistance of Claude (Anthropic). It has not been checked by a human expert.
